\documentclass[12pt,a4paper]{article}
\usepackage[utf8]{inputenc}
\usepackage[romanian]{babel}
\usepackage{geometry}
\geometry{a4paper, margin=2.5cm}
\usepackage{graphicx}
\usepackage{amsmath}
\usepackage{amssymb}
\usepackage{hyperref}
\usepackage{times}
\usepackage{booktabs}
\usepackage{float}
\usepackage{caption}
\usepackage{enumitem}
\renewcommand{\baselinestretch}{1.15}
\begin{document}
% ============================================================
% PAGINA DE TITLU
% ============================================================
\begin{titlepage}
    \centering
    \vspace*{4cm}
    {\Huge \textbf{Adaptarea inteligentă a\\aspect-ratio-ului video fără ML}}\\[1.5cm]
    {\Large Proiect Aplicații Multimedia}\\[2.5cm]
    {\Large \textbf{Autor:} Madar Gabriel}\\[0.4cm]
    {\Large \textbf{Grupa:} 332AA}\\[0.4cm]
    {\Large \textbf{An:} 2026}\\[3cm]
    \vfill
    {\large Universitatea Națională de Știință și Tehnologie\\POLITEHNICA București}
\end{titlepage}
\tableofcontents
\newpage
% ============================================================
% 1. INTRODUCERE
% ============================================================
\section{Introducere, prezentarea temei și a obiectivelor}
Proiectul abordează problema adaptării conținutului video la formate de afișare diferite (retargeting) fără a deforma proporțiile elementelor vizuale importante. Soluția propusă este o implementare tehnică și o analiză de performanță a metodologiei de „Video Seam Carving" bazată pe un model perceptiv JND (Just Noticeable Difference), conform lucrării lui Chiang et al. [2].
Algoritmul funcționează în două moduri complementare: \textbf{contracție} (eliminarea iterativă a cusăturilor de pixeli cu energie minimă, reducând lățimea sau înălțimea cadrului) și \textbf{extindere} (inserarea de noi pixeli prin interpolarea vecinilor de-a lungul acelorași cusături cu energie minimă, mărind astfel dimensiunea cadrului). Ambele moduri utilizează același aparat matematic JND pentru a proteja regiunile vizuale importante.
\subsection{Motivația alegerii temei}
Alegerea temei este motivată de următoarele considerente:
\begin{itemize}[noitemsep]
    \item \textbf{Problema adaptării conținutului multimedia:} Diversitatea dispozitivelor de afișare (monitoare ultra-wide, telefoane verticale, tablete) impune nevoia de a redimensiona inteligent materialele video fără a pierde informația vizuală importantă. Aceasta este o problemă centrală în domeniul aplicațiilor multimedia.
    \item \textbf{Modelarea percepției vizuale umane:} Proiectul oferă oportunitatea de a studia și implementa modele psiho-vizuale (Just Noticeable Difference) care cuantifică matematic limitele sistemului vizual uman, concepte fundamentale în compresia și procesarea multimedia.
    \item \textbf{Procesare video accelerată:} Natura repetitivă a operațiilor per-pixel pe fluxuri video face din accelerarea GPU o necesitate practică, demonstrând cum algoritmii multimedia clasici pot fi adaptați pentru cerințele de performanță ale aplicațiilor moderne.
\end{itemize}
\subsection{Obiectivele proiectului}
\begin{enumerate}[noitemsep]
    \item Implementarea integrală a pipeline-ului de procesare în C++ și CUDA, fără biblioteci pre-procesate pentru algoritmii centrali.
    \item Paralelizarea calculelor de gradient spațial, a filtrelor perceptuale JND și a programării dinamice.
    \item Analiza critică a limitelor abordării deterministe non-ML prin teste de stres riguroase.
    \item Dezvoltarea unei interfețe grafice (GUI) pentru controlul parametrilor și vizualizarea procesului.
\end{enumerate}
% ============================================================
% 2. SUPORT TEHNIC
% ============================================================
\newpage
\section{Suport tehnic, context teoretic și realizări similare}
\subsection{Contextul Multimedia}
Conținutul video este o succesiune temporală de matrice bidimensionale (cadre), fiecare pixel codificând informația de culoare în spațiul RGB sau YUV. Volumul colosal de date necesită algoritmi complecși de compresie: compresia spațială (intra-frame, similară JPEG) și compresia temporală (inter-frame, bazată pe vectori de mișcare, conform standardelor H.264/HEVC).
Conceptul de \textbf{Just Noticeable Difference (JND)}, pragul minim de modificare pe care ochiul uman îl poate detecta, este esențial nu doar în compresia video și watermarking, ci și în algoritmul nostru de retargeting, unde permite ignorarea regiunilor imperceptibile.
\subsection{Trecere în revistă a realizărilor similare}
\begin{itemize}[noitemsep]
    \item \textbf{Metode clasice (Letterboxing, Cropping, Scalare):} Soluții primitive. Letterboxing-ul pierde spațiu util, cropping-ul pierde context, iar scalarea asimetrică distorsionează geometric toate obiectele.
    \item \textbf{Seam Carving (Avidan și Shamir, 2007) [1]:} A introdus eliminarea iterativă a căilor de pixeli cu energie minimă (,,cusături''). Rezultate excelente pe imagini statice. Două probleme majore la video: (a) jitter temporal (cusăturile fluctuează între cadre) și (b) lipsa de protecție semantică (fețele umane pot fi secționate).
    \item \textbf{JND-Based Video Carving (Chiang et al., 2020) [2]:} Hibridizează energia clasică cu modelul psiho-vizual JND, aplatizând energia în regiunile invizibile ochiului. Proiectul de față implementează această metodologie pe arhitectura CUDA.
\end{itemize}
\subsection{Aparatul Matematic}
\subsubsection{Energia Spațială (Gradientul Sobel)}
Energia fundamentală este estimată prin operatorul Sobel:
\begin{equation}
E_{Sobel}(x, y) = |G_{x}(x,y)| + |G_{y}(x,y)|
\end{equation}
unde $G_x$ și $G_y$ reprezintă convoluția imaginii grayscale cu matricele Sobel de $3 \times 3$.
\begin{figure}[H]
\centering
\begin{minipage}{0.48\textwidth}
  \centering
  \includegraphics[width=\linewidth]{mack_original.png}
\end{minipage}\hfill
\begin{minipage}{0.48\textwidth}
  \centering
  \includegraphics[width=\linewidth]{mack_jnd.png}
\end{minipage}
\caption{Caz ideal: Cadrul original (stânga) și harta de energie JND (dreapta). Se observă aplatizarea completă a energiei pe fundalul neted, protejând astfel subiectul central.}
\end{figure}
\subsubsection{Modelul Perceptiv JND Spațial}
Pragul $JND_S$ combină două funcții non-liniare: mascarea luminanței de fundal $f_1(bg)$ și mascarea texturii $f_2(mg)$.
\textbf{Luminanța de fundal} $bg(x,y)$ se calculează ca medie ponderată într-o vecinătate de $5 \times 5$. Funcția $f_1$ reflectă legea lui Weber (sensibilitate ridicată în zone întunecate):
\begin{equation}
f_1(bg) = 
\begin{cases} 
17 \cdot \left(1 - \sqrt{\dfrac{bg}{127}}\right) + 3, & \text{dacă } bg \le 127 \\[6pt]
\dfrac{3}{128} \cdot (bg - 127) + 3, & \text{dacă } bg > 127
\end{cases}
\end{equation}
\textbf{Mascarea texturii} depinde de gradientul maxim calculat pe patru direcții Sobel (0\textdegree, 45\textdegree, 90\textdegree, 135\textdegree), notat $mg(x,y)$. Se definește magnitudinea texturii ca $T_t = \beta \cdot mg$, unde $\beta = 0.15$ este un factor de scalare. Combinarea celor două componente se face printr-o \textbf{selecție condițională}, nu prin înmulțire:
\begin{equation}
JND_S(x,y) = 
\begin{cases}
T_t + JND_{min}, & \text{dacă } mg > T_{prag} \text{ și } T_t > f_1(bg) \\
f_1(bg), & \text{altfel}
\end{cases}
\end{equation}
unde $T_{prag} = 20$ este pragul de activare a mascării de textură, iar $JND_{min} = 3$ este o valoare minimă de prag. Logica reflectă faptul că, într-o zonă puternic texturată, mascarea texturii domină și înlocuiește mascarea de luminanță (dacă o depășește); altfel, pragul rămâne dictat exclusiv de adaptarea la luminanță.
\subsubsection{Componenta Temporală JND}
Pentru fluxuri video, se adaugă componenta temporală. Se calculează mai întâi media diferenței absolute de luminanță (ILD) pe o fereastră de $5 \times 5$ între cadrul curent $I_t$ și cadrul anterior $I_{t-1}$:
\begin{equation}
\overline{ILD}(x,y) = \frac{1}{25} \sum_{dx=-2}^{2} \sum_{dy=-2}^{2} |I_t(x{+}dx, y{+}dy) - I_{t-1}(x{+}dx, y{+}dy)|
\end{equation}
Factorul temporal este apoi:
\begin{equation}
JND_T(x,y) = 1.0 + \alpha_T \cdot \overline{ILD}(x,y)
\end{equation}
unde $\alpha_T = 0.03$. Pragul final combinat este produsul celor două componente: $JND(x,y) = JND_S(x,y) \cdot JND_T(x,y)$.
\textbf{Calculul energiei:} Spre deosebire de o abordare naivă în care s-ar scădea pragul JND din gradientul Sobel, implementarea urmează fidel lucrarea [2]: se \textbf{aplatizează mai întâi intensitățile} pixelilor sub pragul JND (înlocuindu-le cu media locală pe $5 \times 5$), și abia apoi se aplică operatorul Sobel pe imaginea modificată. Această ordine este critică: gradientul rezultat reflectă doar variațiile vizibile ochiului uman, eliminând tranzițiile imperceptibile direct din semnalul de intrare.
\begin{equation}
I'(x,y) = 
\begin{cases}
bg(x,y), & \text{dacă } |I(x,y) - bg(x,y)| < \lambda \cdot JND(x,y) \\
I(x,y), & \text{altfel}
\end{cases}
\end{equation}
\begin{equation}
E(x,y) = 0.7 \cdot \left(|G_x(I')| + |G_y(I')|\right) + 0.3 \cdot \left(|G_x(I)| + |G_y(I)|\right)
\end{equation}
Ponderarea de $70\%/30\%$ asigură că energia finală este dominată de imaginea aplatizată (care ignoră tranzițiile imperceptibile), dar păstrează o componentă din imaginea originală pentru a nu pierde complet muchiile reale șterse accidental de aplatizare.
\subsubsection{Programare Dinamică și Forward Energy}
Costul cumulativ $M(i,j)$ utilizează „Forward Energy" pentru a penaliza muchiile noi create prin eliminarea unui pixel:
\begin{align}
C_L(i, j) &= |I(i, j{+}1) - I(i, j{-}1)| + |I(i{-}1, j) - I(i, j{-}1)| \\
C_U(i, j) &= |I(i, j{+}1) - I(i, j{-}1)| \\
C_R(i, j) &= |I(i, j{+}1) - I(i, j{-}1)| + |I(i{-}1, j) - I(i, j{+}1)|
\end{align}
Recurența:
\begin{equation}
M(i, j) = E(i, j) + \min \left\{
\begin{array}{l}
M(i{-}1, j{-}1) + C_L(i, j) \\
M(i{-}1, j) + C_U(i, j) \\
M(i{-}1, j{+}1) + C_R(i, j)
\end{array} \right.
\end{equation}
% ============================================================
% 3. IMPLEMENTARE CUDA
% ============================================================
\section{Prezentare tehnică a implementării CUDA}
Întreaga logică de procesare a fost implementată de la zero sub forma unor kernel-uri CUDA, conform modelului de programare descris în [3]. OpenCV a fost utilizată exclusiv pentru citirea/scrierea fișierelor video. Arhitectura pipeline-ului este prezentată în Figura~\ref{fig:pipeline}.
\begin{figure}[H]
\centering
\includegraphics[width=\textwidth]{pipeline.png}
\caption{Diagrama pipeline-ului arhitectural CUDA}
\label{fig:pipeline}
\end{figure}
\subsection{Kernel-uri CUDA: structură și optimizări}
În continuare sunt detaliate cele patru etape critice ale pipeline-ului, cu accent pe contribuția personală în optimizarea fiecăreia.
\subsubsection{\texttt{bgrToGrayKernel} - Conversia culorilor}
\begin{itemize}[noitemsep]
    \item \textbf{Intrare:} Buffer BGR interclasat (3 octeți/pixel, layout: B-G-R-B-G-R...).
    \item \textbf{Ieșire:} Buffer Grayscale liniar (1 octet/pixel).
    \item \textbf{Mapare thread-pixel:} Fiecare fir de execuție procesează exact un pixel, calculat prin indexul 2D unic (\texttt{blockIdx}, \texttt{threadIdx}).
    \item \textbf{Formula aplicată:} $L = 0.299 \cdot R + 0.587 \cdot G + 0.114 \cdot B$
    \item \textbf{Optimizare critică:} Firele dintr-un Warp (32 threads) citesc din zone contigue de memorie, asigurând \textbf{Memory Coalescence} și maximizând banda de transfer.
\end{itemize}
\subsubsection{\texttt{spatialJndKernel} - Shared Memory pentru fereastra 5x5}
Calculul hărții JND necesită citirea unei vecinătăți de $5 \times 5$ pixeli pentru fiecare pixel analizat. Într-o abordare naivă, un thread ar efectua 25 de citiri din memoria globală lentă.
\textbf{Strategia de optimizare implementată:}
\begin{enumerate}[noitemsep]
    \item Se alocă \textbf{Shared Memory} de dimensiune $(B_x + 4) \times (B_y + 4)$ pentru un bloc de $B_x \times B_y$ threads. Marginea de 2 pixeli pe fiecare latură acoperă „haloul" necesar ferestrei de $5 \times 5$.
    \item Firele încarcă colaborativ datele proprii și pixelii de graniță (halo pixels) din memoria globală în shared memory.
    \item Bariera \texttt{\_\_syncthreads()} garantează că întreaga fereastră este disponibilă înainte de calcul.
    \item Calculul $f_1(bg)$, $f_2(mg)$ și $JND_S$ se efectuează exclusiv din shared memory, cu latență de ordinul nanosecundelor.
\end{enumerate}
\subsubsection{Programarea Dinamică prin Cooperative Groups}
Recurența DP introduce o \textbf{dependență de date pe verticală}: rândul $y$ nu poate fi procesat până când rândul $y{-}1$ nu este complet finalizat. Abordarea tradițională (720 de lansări de kernel per cadru) introduce overhead masiv pe magistrala PCIe. Soluția adoptată folosește API-ul Cooperative Groups, documentat în [4].
\textbf{Soluția implementată:}
\begin{enumerate}[noitemsep]
    \item Se lansează un \textbf{singur kernel} care acoperă întreaga lățime a cadrului.
    \item Fiecare fir calculează costul pentru pixelul său din rândul curent.
    \item Se invocă bariera globală \texttt{cg::this\_grid().sync()} din API-ul Cooperative Groups.
    \item Toate thread-urile de pe întreg GPU-ul se sincronizează înainte de avansarea la rândul următor.
    \item Rezultat: o singură lansare de kernel per iterație de seam, în loc de $H$ lansări.
\end{enumerate}
\subsubsection{\texttt{removeSeamsPrefixSumKernel} - Warp Shuffle}
După identificarea cusăturii, pixelii marcați trebuie eliminați și memoria compactată. Problema: fiecare pixel supraviețuitor trebuie scris la un index nou, necunoscut a priori.
\textbf{Algoritmul de Prefix Sum (Exclusive Scan):}
\begin{enumerate}[noitemsep]
    \item Fiecare thread generează un flag: \texttt{1} dacă pixelul este păstrat, \texttt{0} dacă este eliminat.
    \item În interiorul unui Warp (32 threads), se aplică instrucțiunea hardware \texttt{\_\_shfl\_up\_sync(mask, val, offset)} cu offset-uri crescătoare ($1, 2, 4, 8, 16$).
    \item După 5 pași (pattern „butterfly"), fiecare thread cunoaște exact suma prefixelor tuturor thread-urilor anterioare.
    \item Această sumă este indexul de scriere unic în buffer-ul de ieșire.
    \item Rezultat: compactare de memorie complet paralelă, fără atomics, fără race conditions.
\end{enumerate}
\subsubsection{\texttt{insertSeamsPrefixSumKernel} - Extinderea cadrului}
Pe lângă eliminarea cusăturilor (contracție), algoritmul suportă și operația inversă: \textbf{extinderea} cadrului prin inserarea de noi pixeli. Procesul se desfășoară în două faze:
\begin{enumerate}[noitemsep]
    \item \textbf{Faza de simulare:} Se identifică toate cele $K$ cusături cu energie minimă prin eliminări succesive pe o copie a imaginii originale, memorând pozițiile acestora într-o hartă cumulativă (\texttt{d\_masterSeamCount}).
    \item \textbf{Faza de inserare:} Pentru fiecare cusătură memorată, se duplică pixelul respectiv prin interpolare cu vecinul adiacent: $p_{nou} = \frac{p_{stânga} + p_{dreapta}}{2}$. \newline
Kernel-ul \texttt{insertSeamsPrefixSumKernel} folosește același algoritm de Prefix Sum bazat pe Warp Shuffle pentru a calcula pozițiile de scriere, însă de data aceasta fiecare pixel generează $1 + \text{count}$ intrări în buffer-ul de ieșire (pixelul original plus pixelii interpolați).
\end{enumerate}
Această abordare garantează că pixelii adăugați sunt plasați exact în zonele cu cea mai mică importanță vizuală, păstrând coerența structurală a cadrului.
\subsection{Optimizarea Multi-Seam (Batching)}
Abordarea clasică cere un ciclu complet (energie $\to$ DP $\to$ eliminare) pentru fiecare pixel de lățime eliminat. Pentru a crește performanța, am implementat un mecanism de \textbf{Batching}: identificarea simultană a $N$ căi de cost minim dintr-o singură hartă energetică, cu o regulă de separare minimă pe axa X.
\begin{table}[H]
\centering
\begin{tabular}{lccc}
\toprule
\textbf{Parametru} & \textbf{Batch = 1} & \textbf{Batch = 10} & \textbf{Observații} \\
\midrule
Calitate vizuală & Maximă & Degradată & Drumuri intersectate \\
Factor de accelerare & $1\times$ & $8{-}9\times$ & Reducere cicluri DP \\
Artefacte geometrice & Minime & Vizibile & Energie statică invalidă \\
\bottomrule
\end{tabular}
\caption{Compromisul calitate vs. performanță în procesarea Multi-Seam}
\end{table}
Degradarea calității apare deoarece eliminarea primei cusături modifică topologia energetică locală. Într-o execuție Batch=10, această modificare este invizibilă algoritmului până la finalizarea lotului, forțând drumurile vecine să traverseze zone care au devenit critice.
% ============================================================
% 4. MOD DE UTILIZARE ȘI REZULTATE
% ============================================================
\section{Mod de utilizare, rezultate și teste de stres}
\subsection{Interfața grafică (GUI)}
Interfața a fost dezvoltată în Python folosind framework-ul NiceGUI, care instanțiază un server web local. Aceasta separă stratul de prezentare (browser) de stratul de procesare (executabil C++ CUDA).
\textbf{Funcționalități principale:}
\begin{itemize}[noitemsep]
    \item Selecția fișierului video și a rezoluției țintă (contracție sau extindere).
    \item Meniu „Advanced Quality Settings" cu parametrii configurabili:
    \begin{itemize}[noitemsep]
        \item \textbf{Lambda:} Intensitatea aplatizării JND (valoare mare = tăiere agresivă în zone netede).
        \item \textbf{Skin Bias:} Penalizare suplimentară pentru pixelii cu profil de culoare al pielii.
        \item \textbf{Motion Weight:} Amplificarea energiei obiectelor în mișcare (protecție temporală).
    \end{itemize}
    \item Modul \textbf{Visualizer}: generează un fișier video intermediar cu traseele de cost minim marcate cu roșu, permițând analiza vizuală a procesului decizional. În modul de extindere, pixelii nou inserați sunt evidențiați distinct.
\end{itemize}
\begin{figure}[H]
\centering
\centering
    \includegraphics[width=1\linewidth]{interfata.png}
\caption{Interfața grafică principală cu meniul de setări avansate}
\end{figure}
\begin{figure}[H]
\centering
 \centering
    \includegraphics[width=1\linewidth]{visualizer.png}
\caption{Modulul Visualizer cu cusăturile marcate cu roșu}
\end{figure}
\subsection{Rezultate de performanță (CUDA)}
Măsurătorile au fost efectuate la parametrul Batch = 1 (calitate maximă, o singură cusătură procesată per iterație):
\begin{table}[H]
\centering
\begin{tabular}{lcc}
\toprule
\textbf{Rezoluție sursă} & \textbf{Operația de retargeting} & \textbf{FPS (Batch=1)} \\
\midrule
$1280 \times 720$ & $1280 \to 780$ lățime (500 cusături) & \textbf{0.9} \\
$640 \times 360$  & $640 \to 440$ lățime (200 cusături) & \textbf{5.0} \\
$360 \times 640$  & $360 \to 260$ lățime (100 cusături) & \textbf{9.0} \\
\bottomrule
\end{tabular}
\caption{Performanța CUDA la Batch = 1 (calitate maximă)}
\end{table}
Aplicând Batch = 10, valorile de mai sus cresc de 8-9 ori, permițând procesare apropiată de timp real la rezoluții SD. Costul este degradarea calității vizuale documentată în Secțiunea 3.2.
\subsection{Limitele algoritmice: eșecuri în teste de stres}
Testarea în condiții limită expune două categorii de eșecuri fundamentale:
\begin{enumerate}
    \item \textbf{Distorsionarea fețelor umane:} Fețele cu piele netedă și iluminată uniform au gradient Sobel scăzut. Algoritmul le consideră „ieftine" și le secționează preferențial față de fundalurile puternic texturate (ziduri de cărămidă, vegetație). Euristicile de Skin Bias atenuează efectul doar parțial.
    
    \item \textbf{Ruperea liniilor drepte:} Elementele rigide rectilinii (stâlpi, cadre de ușă, orizont) sunt treptat dislocate, dobândind un profil de „fierăstrău". Forward Energy previne salturile mari, dar cedează în cele din urmă.
\end{enumerate}
\textbf{Cauza fundamentală:} Analiza bottom-up (per-pixel) nu dispune de înțelegerea semantică necesară pentru a identifica obiecte ca entități coerente. Protecția rămâne o speculație matematică locală.
\begin{figure}[H]
\centering
\includegraphics[width=1\linewidth]{fataTaiata.png}
\centering
    \includegraphics[width=1\linewidth]{cladireStricata.png}
\caption{Artefacte vizuale: distorsionarea fețelor și ruperea liniilor geometrice}
\end{figure}
\begin{figure}[H]
\centering
\includegraphics[width=0.25\textwidth]{qrvideo.png} \\
\vspace{0.2cm}
\href{https://www.youtube.com/watch?v=E84f_W4LxLI}{\textbf{Click aici pentru a viziona demonstrația video}}
\caption{Demonstrativ video al aplicației (Scanare QR sau Click pe link)}
\end{figure}
% ============================================================
% 5. CONCLUZII
% ============================================================
\section{Concluzii}
Obiectivele proiectului au fost îndeplinite:
\begin{itemize}[noitemsep]
    \item Pipeline-ul complet de Video Seam Carving a fost implementat integral pe GPU în CUDA, suportând atât contracția cât și extinderea cadrelor.
    \item Modelul perceptiv JND protejează zonele de interes vizual mai eficient decât un gradient Sobel simplu.
    \item Mecanismul de Batching permite scalarea performanței până la regimul real-time la rezoluții SD.
    \item Testele de stres au documentat limitările inerente ale abordării deterministe.
\end{itemize}
Cu toate acestea, trebuie afirmat critic că soluțiile moderne bazate pe Deep Learning (Vision Transformers, GAN-uri) devansează cu ordine de mărime calitatea vizuală a oricărei metode non-ML. Modelele AI dispun de o înțelegere top-down: știu că „o față umană nu trebuie deformată" independent de valoarea gradientului local. Orice soluție bazată strict pe reguli deterministe bottom-up nu posedă această reprezentare semantică și prezintă limite de rigiditate iremediabile.
Relevanța proiectului rămâne în valoarea sa didactică pentru domeniul aplicațiilor multimedia. Implementarea de la zero a unui pipeline complet de procesare video, de la decodarea cadrelor și calculul gradienților, la modelarea percepției umane prin JND și manipularea adaptivă a conținutului, oferă o înțelegere profundă a principiilor care stau la baza procesării multimedia moderne. Tehnicile de optimizare GPU utilizate (Shared Memory, Warp Shuffle, Cooperative Groups) completează această perspectivă cu abilități practice de procesare paralelă a datelor multimedia la scară largă.
% ============================================================
% 6. NOTĂ LLM
% ============================================================
\section{Notă privind utilizarea instrumentelor de tip Large Language Model}
Conform regulamentului universitar (Regulament UPB), declar că un instrument de tip Large Language Model (LLM) a fost utilizat în următoarele scopuri specifice:
\begin{itemize}[noitemsep]
    \item \textbf{Optimizarea kernel-urilor CUDA:} Rafinarea strategiilor de alocare a Shared Memory și a instrucțiunilor Warp Shuffle (\texttt{\_\_shfl\_up\_sync}) în cadrul kernel-urilor de Prefix Sum.
    \item \textbf{Tehnoredactarea documentului:} Structurarea și formatarea conținutului în format LaTeX, conform cerințelor metodologice impuse.
    \item \textbf{Rafinarea explicațiilor tehnice:} Corectarea gramaticală și formularea coerentă a paragrafelor de analiză tehnică.
\end{itemize}
Întregul aparat matematic, implementarea codului original, măsurătorile empirice de performanță și interpretările critice ale rezultatelor rămân exclusiv produsul analizei și experimentării personale.
% ============================================================
% 7. BIBLIOGRAFIE
% ============================================================
\newpage
\section{Referințe Bibliografice}
\begin{enumerate}
    \item Avidan, S., Shamir, A. (2007). Seam carving for media retargeting. \textit{ACM Transactions on Graphics (TOG)}, 26(3), 10. DOI: 10.1145/1276377.1276390
    \item Chiang, P. P., Wang, C. W., Tsai, C. M., Shen, V. R., Kuo, S. Y. (2020). Fast JND-Based Video Carving with GPU Acceleration for Real-Time Video Retargeting. \textit{IEEE Transactions on Circuits and Systems for Video Technology}, 30(2), 527-540. DOI: 10.1109/TCSVT.2019.2894300
    \item NVIDIA Corporation (Accesat Mai 2026). CUDA C++ Programming Guide. Documentație Internet. \url{https://docs.nvidia.com/cuda/cuda-c-programming-guide/index.html}
    \item NVIDIA Corporation (Accesat Mai 2026). Cooperative Groups: Flexible CUDA Thread Programming. Documentație Internet. \url{https://developer.nvidia.com/blog/cooperative-groups/}
\end{enumerate}
\end{document}
